\begin{definition}
	\textit{Борелевская $\sigma$-алгебра} - это наименьшая (по включению) $\sigma$-алгерба, содержащая все открытые или замкнутые множества в $\R^n$. Все её элементы - это \textit{борелевские} множества.
\end{definition}
\begin{examples}
	$\ds\bigcap_{i=1}^\i G_i$ - открытые, называется \textit{множества типа $G_\delta$}.\\
	$\ds\bc_{i=1}^\i F_i$ - закрытые, называется \textit{множества типа $F_{\sigma}$}
\end{examples}
\begin{note}
	Отметим, типы множеств $\delta$ и $\sigma$ можно использовать последовательно, то есть:
	\[G_{\delta_\sigma}:\bigcap_{i=1}^\i\bc_{j=1}^\i G_{i,j}\]
\end{note}
\begin{example}(Канторово множество)\\
	$I_0=[0,1],\ I_1=I_0\sm J_1=I_{1,1}\cup I_{2,1}$\\
	$I_{1,1}=\l[0,\frac13\r],\ I_{2,1}=\l[\frac23,1\r],\ J_1=\l(\frac13,\frac23\r)$
	$I_2=I_1\sm J_2,\ I_2=I_{1,2}\cup I_{2,2}\cup I_{3,2}\cup I_{4,2},\ J_2=J_{2,1}\cup J_{2,2}$\\
	Канторово множество: $P=\bigcap_{m=1}^\i I_m$.\\
	Заметим, числа из $J_1=\l(\frac13,\frac23\r)$ в троичной системе счисления имеют вид:
	\[(\fa x\in J_1)\ x=0,1\dots_3\]
	Для $J_2=\l(\frac19,\frac29\r)\cup\l(\frac79,\frac89\r)$:
	\[(\fa x\in J_2)\ x=0,01\dots_3\text{ или }x=0,21\dots_3\]
	Таким образом, $P$ состоит из тех и только тех чисел на $[0,1]$, которые могут быть записаны в троичной системе без использования 1.\\
	Построим отображение между $P$ и чем-то: заменим в записи числа из $P$ все 2 на 1, и возьмём не троичную, а двоичную систему.\\
	Получили биекцию между $P$ и $[0,1]$.\\
	Теперь найдём $\mu(P)$:
	\[\mu(P)=\lim_{m\ra\i}|I_m|=1-\sum_{j=1}^\i|J_j|=1-\frac13-2\cdot\frac19-4\cdot\frac1{27}-\dots=1-\frac13\cdot\frac1{1-\frac23}=0\]
	\[\mu^*_J(P)=\mu(P)=0\]
	Теперь будем выкидывать из $[0,1]$ не по $\frac13$ от отрезка, а по $\frac15$. $\t P=\bigcap_{m=0}^\i\t I_m$\\
	\[\mu(\t P)=1-\sum_{j=1}^\i|\t J_j|=1-\frac15-2\cdot\frac1{25}-4\cdot\frac1{125}-\dots=1-\frac15\cdot\frac{1}{1-\frac25}=\frac23\]
	\[\mu_J^*(\t P)=\mu(\t P)=\frac23,\ \mu_*^J(\t P)=0\]
	Заметим: $\fa(a,b)\subset[0,1]\ (a,b)\cap\t P\neq\emptyset\Ra\ex I_{m,k}\cap(a,b)\neq\emptyset$, поэтому внутренняя мера по Жордану 0.
\end{example}
\begin{example}(Неизмеримое по Лебегу множество)\\
	Рассмотрим $[0,1]$.\\
	Введём отношение эквивалентности: 
	\[x\leftrightarrow y\overset{def}\lra x-y\in\Q\] 
	Тогда, по теореме без доказательства, любое множество разбивается на классы эквивалентности: 
	\[[0,1]=\bsc_\al Cl_\al\]
	Теперь воспользуемся "наивной"\ аксиомой выбора.\\
	Пусть $E$ - множество, содержащее ровно по одному элементу из каждого $Cl_\al$\\
	Положим:
	\[\{r_m\}_{m=1}^\i=\Q\cap[-1,1],\ E_m:=E+r_m:=\{x+r_m:x\in E\}\]
	Докажем, что $\fa n,m\in\N\ E_m\cap E_n=\emptyset$.\\
	От противного, пусть $x\in E_m\cap E_n$. Тогда:
	\[x-r_m\in E,\ x-r_n\in E\]
	\[(x-r_m)-(x-r_n)\in\Q\]
	Получили, что $x-r_n$ и $x-r_m$ попадают в один и тот же класс эквивалентности. Противоречие.\\
	Предположим, что $E$ измеримо по Лебегу, $\mu:=\mu(E)$.
	\[(\fa m\in\N)\ \mu(E_m)=\mu,\ \mu\l(\bsc_{m=1}^\i E_m\r)=\sum_{m=1}^\i\mu\]
	\[\bsc_{m=1}^\i E_m\subset[-1,2]\Ra\sum_{m=1}^\i\le3\Ra\mu=0\]
	\[x\in[0,1]\Ra\ex\al:x\in Cl_\al\Ra\ex y\in E\cap Cl_\al: r_m=y-x\in\Q\cap[-1,1]\Ra x\in E_m\]
	\[[0,1]\subset\bsc_{m=1}^\i E_m\Ra\mu([0,1])=1\le\sum_{m=1}^\i\mu\]
\end{example}
\subsection{Измеримые функции}
\begin{definition}
	\textit{Лебеговым множеством} функции $f$, заданной на $E\subset\R^n,\ f:E\ra \R$, называется:
	\[E_f(a)=\{x\in E:f(x)<a\},\ a\in\R\]
\end{definition}
\begin{definition}
	Числовая функция, заданная на измеримом по Лебегу множестве $E\subset\R^n$, называется \textit{измеримой}, если $(\fa a\in\R)\ E_f(a)$ измеримо
\end{definition}
\begin{theorem}
	Совокупность измеримых функций на заданном измеримом множестве $E$ не изменится, если в определении лебегова множества вместо $<$ написать любой из знаков: $\le,>,\ge$.
\end{theorem}
\begin{proof}
	$\ge$ - очевидно.
	\[\{x\in E:f(x)\le a\}=\bigcap_{m=1}^\i E_f\l(a+\frac1m\r)\]
	\[E_f(a)=\{x\in E:f(x)<a\}=\bc_{m=1}^\i\{x\in E:f(x)\le a-\frac1m\}\]
	Получили требуемое.
\end{proof}
\begin{definition}
	\textit{Характеристическая функция} (индикатор) множества $A$:
	\[
	\chi_A(x)=
	\begin{cases}
		1,\ x\in A\\
		0,\ x\notin A
	\end{cases}
	=1_A(x)
	\]
\end{definition}
\begin{note}
	Заметим: $1_A$ измерима $\lra$ A измеримо.
\end{note}
\begin{note}
	Непрерывные функции на $\R^n$ измеримы.
\end{note}
\begin{proof}
	$E_f(a)=f^{-1}((-\i,a))$ - открытое множество.
\end{proof}
\begin{theorem}\nf{(Алгебраическая структура совокупности измеримых функций)}\label{23_th2}
	\begin{enumerate}
		\item Если $f$ непрерывна на $g(E)$, $g$ - измерима на $E\subset\R^n$, то $f\circ g$ измерима.
		\item Если $f,g$ измеримы на $E\subset\R^n$, то $f\pm g,\ f\cdot g$ и $\frac fg$ измеримы на своих областях определения.
	\end{enumerate}
\end{theorem}
\begin{proof}
	\begin{enumerate}
		\item $E_{f\circ g}(a)=\{x\in E:f(g(x))<a\}$, тогда:
		\[\{y\in g(E):f(y)<a\}=f^{-1}((-\i,a))\text{ - открыто в }g(E)\]
		\[(\ex G\text{ - открытое в }\R^n)\ f^{-1}((-\i,a))=G\cap g(E)\]
		\[E_{f\circ g}(a)=g^{-1}(G\cap g(E))=g^{-1}(G)=g^{-1}\l(\bsc_{m=1}^\i P_m\r)=\bc_{m=1}^\i g^{-1}(P_m)\]
		Пусть $P_m=[\al,\beta)$, тогда:
		\[g^{-1}([\al,\beta))=\{x\in E:\al\le g(x)<\beta\}=E_g(\beta)\sm E_g(\al)\]
	\item 
	\begin{multline*}
		E_{f+g}(a)=\{x\in E:f(x)+g(x)<a\}=\{x\in E:f(x)<a-g(x)\}=\\
		\bc_{r\in\Q}\Big(\underbrace{\{x\in E:f(x)\le r\}}_{\text{измеримо}}\cap\underbrace{\{x\in E:g(x)<a-r\}}_{\text{измеримо}}\Big)
	\end{multline*}
	Для разности аналогично.
	\[
	E_{cf}(a)=\{x\in E:cf(x)<a\}=
	\begin{cases}
		\{x\in E:f(x)<\frac ac\}, c>0\\
		\emptyset,\ c=0,\ a\le 0\\
		E,\ c=0,\ a>0\\
		\{x\in E:f(x)>\frac ac\},\ c<0
	\end{cases}
	\]
	Все из этих вариантов измеримы, значит $E_{cf}$ измеримо.
	\[f\cdot g=\frac14((f+g)^2-(f-g)^2)\]
	\[E_{\frac1g}(a)=\{x\in E:\frac1{g(x)}<a\}=
	\begin{cases}
		\{x\in E:g(x)<0\}\cup\{x\in E:g(x)>\frac1a\},\ a>0\\
		\{x\in E:g(x)<0\},\ a=0\\
		\{x\in E:0>g(x)>\frac1a\},\ a<0
	\end{cases}
	\]
	\end{enumerate}
	Получили требуемое.
\end{proof}
\begin{theorem}\nf{(Предельный переход и измеримость функций)}\label{23_th3}\\
	Если последовательность измеримых на множестве $E$ функций $\{f_m\}$ сходится поточечно на $E$ к $f$, то $f$ измеримо.
\end{theorem}
\begin{proof}
	\[\lim_{m\ra\i}f_m(x)=\vlim[m\ra\i]f_m(x)=\lim_{m\ra\i}\sup_{k\ge m}f_k(x)=\inf_{m\in\N}\sup_{k\ge m}f_k(x)\]
	Докажем, что $\ds\sup f_k$ измерим.
	\[\{x\in E:\sup f_k(X)\le a\}=\bigcap_{k}\{x\in E:f_k(x)\le a\}\]
	$\inf$ доказывается аналогично, значит получили требуемое.
\end{proof}
\begin{note}
	Все сформулированные результаты справедливы и для функций, принимающих значения $+\i$ и (или) $-\i$ на измеримых множествах.
\end{note}